\section{敏感性分析与稳健性检验}

由于本文研究的布图规划问题属于离散几何优化，轮廓边长、模块位置与旋转状态的变化通常引起指标的离散响应。因此，本节不对整数决策变量作微分意义下的局部敏感度分析，而是围绕轮廓边长扰动、Terminal 边界口径和跨实例结果一致性进行检验，以判断主要结论是否依赖于个别参数或单一算例。

\subsection{敏感性分析}

\subsubsection{轮廓边长扰动}

由式~\eqref{eq:q3_deadspace}可知，在硬模块总面积 $A_{\mathrm{block}}$ 不变时，死区率 $\Gamma(S)$ 随正方形轮廓边长 $S$ 严格单调增加。以问题三获得的当前最小成功构造边长 $S_{\mathrm{found}}$ 为中心，分别计算 $S_{\mathrm{found}}-1$、$S_{\mathrm{found}}$ 与 $S_{\mathrm{found}}+1$ 对应的理论死区率，结果见表~\ref{tab:sensitivity_side}。

\begin{table}[H]
  \centering
  \caption{轮廓边长扰动下的理论死区率}
  \label{tab:sensitivity_side}
  \small
  \renewcommand{\arraystretch}{1.15}
  \begin{tabular}{c r r r}
    \toprule
    实例 & $\Gamma(S_{\mathrm{found}}-1)$/\% & $\Gamma(S_{\mathrm{found}})$/\% & $\Gamma(S_{\mathrm{found}}+1)$/\% \\
    \midrule
    \texttt{n100} & 6.8763 & 7.3649 & 7.8546 \\
    \texttt{n200} & 5.7286 & 6.2198 & 6.7122 \\
    \texttt{n300} & 5.1711 & 5.5639 & 5.9575 \\
    \bottomrule
  \end{tabular}
\end{table}

表中 $S_{\mathrm{found}}-1$ 对应的数值仅用于衡量目标函数的变化幅度，并不表明该边长下必然存在可行布局。边长增加 $1$ 个单位时，三组实例的死区率分别增加约 $0.49$、$0.49$ 和 $0.39$ 个百分点，变化方向与理论单调性完全一致，不会因数值舍入造成相邻整数边长的目标排序反转。

三组实例的理论搜索下界为 $424$、$420$、$523$，而 $S_{\mathrm{found}}$ 为 $439$、$432$、$537$，分别高出 $15$、$12$、$14$ 个单位，相对差距仅为 $3.54\%$、$2.86\%$、$2.68\%$。这表明当前方案均已进入面积下界附近的紧轮廓区域；但 MaxRects 仍是启发式构造器，上述差距不能作为 $S_{\mathrm{found}}$ 全局最优的证明。

\subsubsection{Terminal 边界口径}

正文将芯片轮廓解释为硬模块的摆放边界，固定 Terminal 仅参与 HPWL 计算。若采用更严格的替代口径，要求所有 Terminal 也位于正方形轮廓内，则根据 Terminal 坐标范围，\texttt{n100}、\texttt{n200} 和 \texttt{n300} 的边长至少需要增至 $444$、$438$ 和 $548$。相对主模型边长分别增加 $5$、$6$ 和 $11$ 个单位，对应死区率为
\begin{equation}
  9.8245\%,\qquad 9.1909\%,\qquad 9.9330\%.
  \label{eq:terminal_sensitivity}
\end{equation}
较主模型结果分别上升 $2.46$、$2.97$ 和 $4.37$ 个百分点。因此，Terminal 是否受轮廓约束会对问题三的死区率数值产生明显影响，是本题中需要明确的重要边界口径。

即便采用该替代口径，三组死区率仍均低于问题二设定的 $15\%$，因而“收紧轮廓能够提高面积利用率”的主要结论不受题意解释改变。正文保留硬模块轮廓口径，并将替代口径的诊断计数列于附录 A.2。

\subsection{稳健性检验}

\subsubsection{几何可行性与指标复算}

对问题一至问题三的最终布局，程序独立检查模块名称完整性、轮廓边界、两两非重叠、模块总面积和 HPWL 复算一致性。三组实例均未出现模块缺失、重复、越界或重叠，面积与线长复算结果也与正文表格一致。因此，尽管大规模组合优化采用启发式算法，不具有全局最优性保证，所得坐标方案仍可确认为满足原模型几何约束与指标口径的合法可行解。

问题四则具有更强的最优性证据：四个模块的实际面积之和为 $24$，任意可行外接轮廓的面积不小于 $24$；完整有限枚举又成功构造出面积恰为 $24$ 的 $6\times4$ 布局。因此，该结果由理论下界与可行构造共同证明为全局最优。

\subsubsection{跨实例变化规律}

在问题二中，局部重插入使 \texttt{n100}、\texttt{n200} 和 \texttt{n300} 的 HPWL 分别降低 $1.7809\%$、$0.3709\%$ 和 $0.6185\%$。三组实例的改善幅度不同，但变化方向一致，说明“网络驱动目标位置---几何合法化---局部重插入”的分层流程能够在不破坏几何可行性的前提下，稳定回收合法化阶段造成的部分线长损失。

在问题三中，三组实例的死区率均由 $15\%$ 降至 $5.5639\%$--$7.3649\%$，同时 HPWL 较问题二分别增加 $8.7489\%$、$7.7917\%$ 和 $14.3094\%$。不同规模网表均呈现“轮廓收紧---线长增加”的同向规律，说明该权衡并非单一算例的偶然结果，而与紧轮廓减少模块位置调整空间的机制相一致。

综合上述检验，本文方案在几何约束、指标复算与跨实例变化规律三个层面均具有一致性。Terminal 边界口径虽会改变问题三的具体死区率，但不会改变轮廓收紧能够提高面积利用率，以及面积利用率与互连线长之间存在权衡的核心结论。
